\section{EconML reuse: preserve the target, history and sampling experiment}
\label{exec19:econml-scope}

A bank choosing which customers should receive a servicing offer needs an
effect for a specified population and decision time. A flexible treatment-effect
package can supply useful estimation machinery, but the meaning of the output
depends on the effect model, the variables used for adjustment and the population
over which predictions are averaged. A repeated-contact application adds a
second obligation: the computational history must preserve the customer's
actual sequence of decisions. These obligations motivate a selective EconML
reference profile, with a separate contract for a future executable adapter.

The primary comparison uses the Python repository \texttt{py-why/EconML} at
commit \texttt{d625dd696ea28c62e6e4d8baea6e4154b98c30b4}, dated 1 October
2026. The source declares version 0.17.0, and its README records that version's
31 July 2026 release. Seventeen selected original files cover package notices,
dependency declarations, static DML, DynamicDML, common inference machinery
and selected upstream tests. The immutable source revision is the identity
used here; it does not authenticate every wheel or installation labelled
0.17.0. The selected files support source-level conclusions, while an actual
native fit and complete estimator-family qualification remain separate work.
\citep{exec19_econml_source}

\subsection{A residual projection can answer a different banking question}
\label{exec19:econml-target}

Write $A$ for an offer indicator, $Y$ for its monetary outcome, $X$ for
pre-action characteristics used to describe effect heterogeneity, and $W$ for
additional adjustment variables. In the selected DML formulation, the nuisance
functions are $m(X,W)=E[A\mid X,W]$ and $q(X,W)=E[Y\mid X,W]$.
Residuals $\widetilde A=A-m(X,W)$ and
$\widetilde Y=Y-q(X,W)$ enter the objective
\begin{equation}
 \widehat\theta\in\mathop{\rm arg\,min}_{\theta\in\Theta}
 \frac{1}{n}\sum_{i=1}^n
 \{\widetilde Y_i-\theta(X_i)\widetilde A_i\}^{2}.
 \label{exec19:econml-objective}
\end{equation}
The selected documentation states a structural residual relationship and
orthogonality conditions; the R-learner implementation forms the corresponding
residuals and final regression. A causal reading additionally needs the stated
identifying information and an appropriate effect-function class $\Theta$.
An accurate nuisance model alone cannot correct a deliberately inadequate
effect model. Variables admitted to $W$ also need causal justification:
automatic adjustment for a post-offer variable can change the question.
\citep{exec19_econml_source}

For a binary offer under a sufficient observed adjustment variable $S$, let
$e(S)=P(A=1\mid S)$ and $\tau(S)=E[Y(1)-Y(0)\mid S]$. With exact
nuisances and a constant final effect, differentiating the population squared
loss gives
\begin{equation}
 \theta_*=
 \frac{E[\widetilde A\widetilde Y]}{E[\widetilde A^2]}
 =\frac{E[e(S)\{1-e(S)\}\tau(S)]}
        {E[e(S)\{1-e(S)\}]}.
 \label{exec19:econml-projection}
\end{equation}
The second equality uses the identified binary-treatment conditional mean
relationship and a positive residual variance. Thus the constant projection
weights strata by treatment variation. The population average effect instead
uses the target population's stratum probabilities. They agree under a
correct constant-effect model, or particular balancing conditions, but need
not agree when effect heterogeneity has been omitted.
\citep{exec19_econml_reference}

Consider a fully specified example with $S\sim\mathrm{Bernoulli}(1/2)$,
$e(0)=1/2$, $e(1)=1/10$, and potential outcomes $Y(a)=100Sa$ in USD.
The offer is assigned using only observed $S$, so there is no hidden
confounding in this construction. The two strata have effects USD 0 and
USD 100. Direct enumeration gives
\begin{equation}
 E[Y(1)-Y(0)]=50,\qquad
 E[\widetilde A\widetilde Y]=\frac92,\qquad
 E[\widetilde A^2]=\frac{17}{100},\qquad
 \theta_*=\frac{450}{17}\simeq26.47.
 \label{exec19:econml-money}
\end{equation}
At an assumed offer cost of USD 35, offering to the whole target population
has expected net value USD 15 per customer. Labelling the constant projection
as that population effect would instead report net value
$450/17-35=-145/17$, about USD $-8.53$. This is a consequence of the
specified misspecification and target mismatch, not a demonstrated EconML
error under a correctly specified effect model. Including the effect modifier
in an adequate $\Theta$ and then averaging over the declared target population
is one route to the intended question in this example.
\citep{exec19_econml_reference}

The base \texttt{ate} method averages predicted contrasts over the supplied
$X$ rows. It cannot recover missing heterogeneity merely by averaging the
same fitted constant repeatedly. A proposed adapter must therefore record
the contrast, treatment scale and reference category, effect-function class,
prediction population and aggregation weights. A marginal effect per unit
of a continuous dose also needs conversion before being interpreted as the
effect of a finite banking action. These fields are part of the scientific
request, not incidental serialization metadata.~\citep{exec19_econml_source}

\subsection{Cross-fitting, numerical conventions and uncertainty have different roles}
\label{exec19:econml-inference}

The inspected final wrapper offers a feature--treatment-residual linear path
and an alternative weighted path. The latter divides the outcome residual
by a signed treatment residual whose magnitude is clipped below at $10^{-5}$,
and uses the original squared treatment residual as a weight. That convention
needs to be preserved in a numerical comparison; it does not supply missing
treatment support. The wrapper should report the convention and diagnostics
without presenting a finite result as evidence that every requested contrast
has adequate overlap.~\citep{exec19_econml_source}

The common learner uses group-aware folds when groups are supplied and a
compatible splitter is used; setting \texttt{cv=1} explicitly disables
cross-fitting. The selected static linear covariance implementation offers
HC0, HC1 and nonrobust choices. It uses a pseudoinverse in the variance
calculation and warns that inference is invalid when the design is rank
deficient. Keeping a borrower's rows in one nuisance fold is useful for
preventing a particular form of training leakage. It does not itself change
the covariance formula into a cluster or calendar estimator. Rank and
sampling warnings must therefore survive the proposed CLI or MCP boundary,
including when the backend returns numerical arrays.
\citep{exec19_econml_source}

A small construction makes the sampling issue explicit. Suppose four customer
score contributions all equal a shared fair shock $Z\in\{-1,1\}$. Each
has mean zero and variance one, but their average equals $Z$ and hence has
variance one. A working calculation that treats the four contributions as
independent gives $1/4$. Group labels in a fitting call cannot remove the
off-diagonal covariance terms. This example concerns a declared common-shock
law; it is not a measured coverage result for EconML.
\citep{exec19_econml_reference}

\subsection{A dynamic panel requires an explicit clock and complete histories}
\label{exec19:econml-panel}

The DynamicDML guide considers treatments assigned over time with all relevant
dynamic confounders observed. Its displayed structural family combines a
state transition linear in previous treatment and state, a potentially
nonlinear assignment policy, and an outcome equation with baseline effect
heterogeneity. The selected implementation targets the final-period outcome,
fits treatment effects backwards through the periods, and uses the first
period's $X$ for heterogeneity. Later $X$ enters as time-varying controls.
These are meaningful design choices for a specified sequential problem; they
do not establish an unrestricted hidden-state servicing or optimal-policy
value guarantee.~\citep{exec19_econml_source}

The code assigns a group's period from its occurrence count in the input.
Its prefit check requires each group to have the same number of rows, but
equal counts do not verify timestamps. For example, the ordered rows
$(a,0),(b,1),(a,1),(b,0)$ contain two dates for each customer. Occurrence
counting nevertheless assigns customer $b$'s date 1 to period 0 and date 0
to period 1. A future adapter must require unique unit--time keys and
canonical chronological histories before making the call. It should also
verify identical unit alignment across period slices, because period-specific
arrays are subsequently combined positionally.
\citep{exec19_econml_source,exec19_econml_reference}

An unbalanced panel raises an error whose text suggests padding the data.
Padding has scientific consequences: an absent contact, an unobserved
repayment and an observed zero are different events. The proposal therefore
requires a declared observation and censoring treatment before any such
conversion. The same principle applies to weights. Although the selected
\texttt{fit} signature accepts sample weights and sample variances, its code
warns that these are unsupported and passes neither into fitting. A requested
weighted population answer must be refused under this profile, rather than
returned as an unweighted result with its warning discarded.
\citep{exec19_econml_source}

The dynamic analytic covariance does include cross-period within-unit score
products. It constructs a block sensitivity matrix and a block score
covariance, then divides the resulting sandwich by the number of groups.
This deserves recognition as more than a rowwise static variance calculation.
It still does not establish covariance qualification for shared calendar
shocks across customers. The proposed request must declare both the history
model and the between-unit sampling law.~\citep{exec19_econml_source}

\subsection{The inspected bootstrap composition can change a history}
\label{exec19:econml-bootstrap}

The source exposes an additional issue that should be resolved before dynamic
bootstrap admission. The common CATE inference mapping selects
\texttt{BootstrapInference}; that object calls \texttt{BootstrapEstimator.fit}.
The latter asks the estimator for strata. DynamicDML returns its group labels.
The bootstrap then draws rows with replacement inside each distinct stratum
and concatenates the draws. Consequently, this composition holds the original
units fixed and resamples their time rows internally. It is not a draw of
whole ordered histories from a population of units. This conclusion follows
from the inspected source call path; no installed backend was run here.
\citep{exec19_econml_source}

For an authored two-period history $(0,1)$, two independent uniform draws from
its rows give
\begin{equation}
 (0,0),\ (0,1),\ (1,0),\ (1,1),
 \qquad P(\hbox{original ordered history})=\frac14.
 \label{exec19:econml-resampling}
\end{equation}
Only $(0,1)$ preserves that history; for two independent unit draws the
probability that both histories survive is $1/16$. Since the selected dynamic
code infers time from occurrence order, duplicated or reversed rows can be
treated as successive periods. This enumeration establishes a contract risk,
not a measured coverage shortfall or a general indictment of the package's
other inference modes.~\citep{exec19_econml_reference}

The selected upstream dynamic tests exercise shapes and inference choices,
including a two-draw bootstrap configuration, and compare selected panel
layouts under automatic analytic inference. Those layouts preserve the
intended chronological information. The selected common bootstrap tests
exercise interval shapes, options and treatment-stratified examples. These
tests were inspected rather than executed. They do not supply the proposed
whole-history conformance evidence. The next prototype should instrument the
pinned resampler with unit, row and time identifiers, preserve the observed
draws, and compare the actual path against an independently specified panel
sampling experiment. A whole-unit bootstrap would itself need an appropriate
independent-unit law; common calendar dependence cannot be repaired by that
substitution without a separate argument.~\citep{exec19_econml_source}

\subsection{Selective reference now, a testable adapter proposal next}
\label{exec19:econml-route}

EconML is a useful reference for flexible effect-function estimation and for
comparison with separately qualified DoubleML profiles. The proposed route
is a narrow Python boundary that preserves targets, histories, warnings and
inference choices. Direct adoption of every estimator and a wholesale port
would enlarge the verification obligation without resolving the scientific
questions above. A transparent finite residual-moment implementation provides
a cheap numerical reference; an estimator with a different longitudinal
contract may be preferable when the selected dynamic structural family is
inappropriate. Refusal remains necessary when identification is unsupported.

The package's root notice contains MIT terms and explicit BSD 3-Clause notices
for specified scikit-learn-derived tree and generalized-random-forest code.
The pinned metadata declares Python at least 3.9, eleven direct dependency
bounds and optional integrations. Its native build requirements include Cython
at least 3.0.11 and below 3.3. These declarations inform feasibility and
maintenance costs; they are not a resolved environment, a complete transitive
notice inventory or proof that a future installation reproduces the selected
source. A prototype needs exact artifacts and applicable notices, with optional
services enabled only when that profile requires them.
\citep{exec19_econml_source}

The first public or synthetic prototype should keep one identified static
target, an explicit effect-function class and a fixed prediction population.
It should reproduce the exact nuisance oracle, test keyed reordering and
treatment-scale conversions, and retain the misspecified constant comparison.
A subsequent dynamic prototype should add complete-history and weight
refusals, trace resampling, and compare declared independent-unit and
common-shock designs. Coverage, interval availability, width and total
failed-run cost should be specified before final evaluation. Such work can
establish whether a narrow implementation merits an expensive bank trial.
The assumed USD amounts above motivate the decision; observed bank prices,
assignment truth and institutional acceptance remain later evidence.
\citep{exec19_econml_reference}
